% Generated draft: requires mathematical and source review.
\documentclass[11pt]{article}
\usepackage[margin=25mm]{geometry}
\usepackage{fontspec}
\setmainfont{Latin Modern Roman}
\usepackage{amsmath,amssymb,mathtools,unicode-math}
\setmathfont{Latin Modern Math}
\usepackage{enumitem,needspace,xurl,hyperref}
\setlength{\parindent}{0pt}
\setlength{\parskip}{5pt}
\setlength{\emergencystretch}{3em}
\newcommand{\N}{\mathbb{N}}
\newcommand{\Z}{\mathbb{Z}}
\newcommand{\Q}{\mathbb{Q}}
\newcommand{\R}{\mathbb{R}}
\newcommand{\C}{\mathbb{C}}
\newcommand{\F}{\mathbb{F}}
\newcommand{\A}{\mathbb{A}}
\newcommand{\PP}{\mathbb{P}}
\newcommand{\E}{\mathbb{E}}
\newcommand{\eps}{\varepsilon}
\newcommand{\dd}{\,\mathrm{d}}
\begin{document}
\section*{UnsolvedMath ID: 10600033}
\textbf{Source problem code:} catalogue entry
\subsection{Definitions and mathematical statement}
Let \(H\) be a finite-dimensional Hopf algebra over a field \(k\). A right integral functional may be specified by a nonzero \(\lambda\in H^*\) satisfying \(\lambda*f=f(1)\lambda\) for every \(f\in H^*\), with convolution \((\lambda*f)(h)=\lambda(h_{(1)})f(h_{(2)})\). Equivalently, in this convention, \((\lambda\otimes\operatorname{id})\Delta(h)=\lambda(h)1\). A quasitriangular structure supplies an invertible universal \(R\)-matrix; any further ribbon, framing and normalization data required by a KRH construction must be stated separately.

The source asks two things:
\begin{enumerate}
\item Give an effective method for computing the Kauffman--Radford--Hennings (KRH) invariants of \(3\)-manifolds using right integrals on finite-dimensional Hopf algebras.
\item Categorify the KRH invariants for \emph{all} finite-dimensional quasitriangular Hopf algebras: construct a suitable categorical or homological invariant whose specified decategorification recovers the intended KRH invariant.
\end{enumerate}
The excerpt gives no computational input representation, complexity bound, categorification category or decategorification map. Nor does it give the additional algebraic hypotheses and nonvanishing normalization conditions needed to define the invariant in the stated full generality; the word all is preserved as a requirement to address, rather than replaced by a narrower known class.

The source explains that KRH invariants are defined on rotational virtual links and, after normalization, are invariant under Kirby-calculus moves, giving invariants of ordinary and virtual \(3\)-manifolds. In the usual rotational convention, the virtual first Reidemeister move is excluded; framing and the remaining move conventions still need specification. It presents rotational virtual knot theory as the proper domain for quantum link invariants generally and views KRH as a test case. This broad contextual viewpoint is retained as the source's motivation, not an independently verified universal theorem.
\subsection{Short English statement}
Make KRH invariants effectively computable using Hopf-algebra right integrals, and categorify them across the full stated class of finite-dimensional quasitriangular Hopf algebras, including the needed normalization data.
\subsection{Sources}
\begin{itemize}
\item \emph{Unsolved Problems in Virtual Knot Theory and Combinatorial Knot Theory}, revised survey, problem-list item 33. \url{https://arxiv.org/abs/1409.2823}. Recorded access 2026-07-29; extraction source-tex. The invariant is attributed by name to Kauffman, Radford and Hennings; its inline reference is absent.
\item Supplied summary references, all incomplete: \emph{Centrality and the KRH Invariant} (2021, arXiv:2107.xxxx); \emph{Quantum Invariants of Links and 3-Manifolds with Boundary defined via Virtual Links} (2021, arXiv:2108.xxxx), with examples (2022, arXiv:2203.xxxx); \emph{Logarithmic Hennings invariants for restricted quantum sl(2)} (2017, arXiv:1705.xxxx); \emph{Integrality and Gauge Dependence of Hennings TQFTs} (2013, arXiv:1305.xxxx). The xxxx strings are placeholders, not complete arXiv identifiers.
\item Local UnsolvedMath registry, \texttt{data/unsolvedmath/problems.json}, permanent ID \(10600033\), exact source code \texttt{catalogue entry}; supplied partially\_solved/PARTIAL-PROGRESS status.
\end{itemize}
\end{document}
